
\chapter{فلسفه توسعه و روش یادگیری}

\section{اصل اول: پروژه باید قابل فهم بماند}

هر کتابخانه یا Framework هزینه ذهنی دارد. اضافه‌کردن Dependency زمانی مجاز است که:
\begin{enumerate}
  \item یک نیاز واقعی را حل کند؛
  \item راه ساده‌تر built-in کافی نباشد؛
  \item Maintenance مناسبی داشته باشد؛
  \item بتوانیم توضیح دهیم چه کاری انجام می‌دهد؛
  \item Test و Debug آن قابل مدیریت باشد.
\end{enumerate}

\section{اصل دوم: Learn Before Magic}

\begin{learningbox}
قرار نیست Code Generation جای یادگیری را بگیرد. هنگام پیاده‌سازی هر بخش ابتدا Concept توضیح داده می‌شود، سپس مثال کوچک، بعد کد واقعی InsightInk.
\end{learningbox}

مثال:
\begin{enumerate}
  \item ابتدا مفهوم Class در Python؛
  \item بعد مفهوم Django Model؛
  \item بعد رابطه Model با Table؛
  \item سپس ساخت \eng{Idea}؛
  \item سپس Migration؛
  \item سپس Query با ORM؛
  \item سپس Test.
\end{enumerate}

\section{اصل سوم: از دانسته‌های Node.js استفاده می‌کنیم}

توسعه‌دهنده با \eng{Node.js/Express} آشناست. بنابراین برای کاهش بار یادگیری، Concepts جدید با مفاهیم آشنا مقایسه می‌شوند:

\begin{longtable}{p{5.6cm}p{6.5cm}}
\toprule
\textbf{Node/Express} & \textbf{Python/Django} \\
\midrule
\endhead
\eng{npm package} & Python package / pip dependency \\
\eng{node\_modules} & virtual environment site-packages \\
\eng{Express Router} & Django URLConf + Views \\
\eng{Middleware} & Django Middleware \\
\eng{Prisma/Sequelize Model} & Django Model/ORM \\
\eng{Migration tool} & Django Migrations \\
\eng{req/res} & Django request/response \\
\eng{JWT/session middleware} & Django authentication/session/token strategy \\
\eng{validation schema} & DRF Serializer validation \\
\bottomrule
\end{longtable}

\section{اصل چهارم: کوچک‌ترین سیستم درست}

در هر فاز ابتدا ساده‌ترین نسخه صحیح ساخته می‌شود.

مثلاً Recommendation:
\begin{enumerate}
  \item Interest match؛
  \item Save/Spark signal؛
  \item Follow signal؛
  \item Jaccard related ideas؛
  \item فقط اگر لازم شد TF-IDF.
\end{enumerate}

نه اینکه از روز اول ML pipeline ساخته شود.

\section{اصل پنجم: Database Source of Truth}

تمام داده‌های ماندگار باید در PostgreSQL باشند. Redis، WebSocket یا Browser cache جایگزین Source of Truth نیستند.

\section{اصل ششم: هر Feature شش حالت دارد}

برای هر Screen/Feature حداقل این Stateها بررسی می‌شوند:
\begin{itemize}
  \item Idle؛
  \item Loading؛
  \item Success؛
  \item Empty؛
  \item Error؛
  \item Offline/Degraded در صورت مرتبط بودن.
\end{itemize}

\section{اصل هفتم: Definition of Done واقعی}

Feature فقط وقتی Done است که:
\begin{itemize}
  \item رفتار اصلی کار کند؛
  \item Validation داشته باشد؛
  \item Permission بررسی شده باشد؛
  \item Error state وجود داشته باشد؛
  \item ترجمه فارسی/انگلیسی در نظر گرفته شده باشد؛
  \item Test مرتبط پاس شود؛
  \item Documentation در صورت تغییر Contract به‌روز شود.
\end{itemize}

\section{اصل هشتم: هر Commit یک مفهوم}

Commit باید کوچک و قابل مرور باشد. نمونه:
\begin{latin}
\begin{lstlisting}
feat(content): add Idea and Source models
test(auth): cover moderator permission boundaries
docs(api): document save endpoint
fix(i18n): correct RTL spacing in idea cards
\end{lstlisting}
\end{latin}

\section{اصل نهم: Performance قبل از Optimization کور}

تا وقتی اندازه‌گیری نکرده‌ایم، Optimization پیچیده نمی‌کنیم. Index، Cache، Redis و Queue بعد از مشاهده نیاز واقعی وارد می‌شوند.

\section{اصل دهم: Project Journal}

در هر فاز بهتر است یک فایل کوتاه Log نگهداری شود که پاسخ دهد:
\begin{itemize}
  \item امروز چه چیزی یاد گرفتم؟
  \item چه چیزی ساختم؟
  \item چه Bugی دیدم؟
  \item چه Decisionی گرفتم؟
  \item مرحله بعد چیست؟
\end{itemize}

این Log در دفاع و گزارش نهایی ارزش زیادی دارد.
